\section{Controlled Case Study}
\label{sec:case}

The case study asks whether capability gaps identified by AUF manifest as observable
grounding and validation failures. It uses an authored two-line bottling plant (17 assets,
73 signals, 22 shared descriptions, metric and imperial units, four interlocks, 14 feed
relations) with 35 tasks that declare their required facts and 137 proposed actions
(84 injected faults in 12 families, 53 legal). Context conditions C0--C4 expose
progressively more capabilities: tags and values; descriptions and unit symbols; a typed
\opcua{}/AAS-like profile; units, topology and taxonomy; and behaviour, guards and
provenance. Validator tiers G0--G4 check the same capability sets. Sufficiency is the
fraction of tasks whose declared facts are all present within a 6,000-token budget
(deterministic surrogate tokenizer); it is not a bound on model accuracy. Tasks, faults and
validators share one authored specification, so the study illustrates mechanisms and is
not an external validation.

\begin{table*}[t]
\centering\footnotesize
\setlength{\tabcolsep}{3pt}
\caption{CS1: progressive capability exposure on the authored reference plant (17 assets, 73 signals, 35 tasks, 84 faulty and 53 legal proposals). Suff.: declared required facts present at 6,000 surrogate tokens (not an accuracy bound). Context C$k$ is paired with validator G$k$.}
\label{tab:cs1}
\begin{tabular}{@{}llrrrlrr@{}}
\toprule
Ctx. & Projection & Top-1 & Verif. & Suff. & Val. & Recall & FPR \\
\midrule
C0 & OPC DA / PLC symbol table & 0\% & 0\% & 0\% & G0 & 0\% & 0\% \\
C1 & SCADA/historian CSV & 79\% & 0\% & 0\% & G1 & 14\% & 0\% \\
C2 & OPC UA NodeSet2 / AAS submodels & 95\% & 95\% & 23\% & G2 & 68\% & 8\% \\
C3 & RDF/OWL + QUDT + Brick + ISA-95 & 100\% & 100\% & 51\% & G3 & 74\% & 0\% \\
C4 & PackML + interlocks + SOSA/PROV & 100\% & 84\% & 86\% & G4 & 100\% & 0\% \\
\bottomrule
\end{tabular}
\end{table*}


\paragraph{CS1: Progressive exposure.} Top-1 grounding rises from \CsTopCOne{} with
descriptions to \CsTopCTwo{} with a typed profile, yet sufficiency at C2 is only
\CsSuffCTwo{}, rising to \CsSuffCFour{} at C4 (Table~\ref{tab:cs1}). Validator recall rises
from \CsRecallGTwo{} at G2, which also raises false alarms on \CsFPRGTwo{} of legal actions
through literal unit comparison, to \CsRecallGFour{} at G4 without false alarms. Typing
resolves the target but not whether the action is admissible.

\begin{table}[t]
\centering\footnotesize
\setlength{\tabcolsep}{3pt}
\caption{CS2: retrieval strategy crossed with exposed capabilities at 6,000 surrogate tokens. Membership and graph retrieval share a fixed entity resolver; graph traversal requires C3/C4 relations.}
\label{tab:cs2}
\begin{tabular}{@{}llrrr@{}}
\toprule
Ctx. & Retrieval & Suff. & Coverage & Tokens \\
\midrule
C2 & lexical & 23\% & 55\% & 5193 \\
C2 & membership & 23\% & 48\% & 1906 \\
C3 & lexical & 40\% & 67\% & 5075 \\
C3 & membership & 29\% & 59\% & 2507 \\
C3 & graph & 46\% & 76\% & 5341 \\
C4 & lexical & 71\% & 85\% & 5107 \\
C4 & membership & 46\% & 69\% & 2767 \\
C4 & graph & 80\% & 93\% & 5345 \\
\bottomrule
\end{tabular}
\end{table}

\begin{table}[t]
\centering\footnotesize
\setlength{\tabcolsep}{3pt}
\caption{CS3: removing one AUF capability at a time from the full C4/G4 configuration. Abstentions on missing evidence are not counted as detections.}
\label{tab:cs3}
\begin{tabular}{@{}lrrr@{}}
\toprule
Removed & Suff. & Recall & Abstain \\
\midrule
none & 80\% & 100\% & 0 \\
behaviour & 63\% & 89\% & 37 \\
constraints & 69\% & 96\% & 20 \\
provenance & 74\% & 85\% & 40 \\
topology & 63\% & 94\% & 5 \\
units & 80\% & 61\% & 60 \\
\bottomrule
\end{tabular}
\end{table}


\paragraph{CS2: Retrieval.} At equal budget, graph traversal yields the highest
sufficiency where relations exist, membership-scoped retrieval is cheapest but omits
evidence, and lexical retrieval cannot compensate for capabilities the context lacks
(Table~\ref{tab:cs2}). Evidence availability and retrieval are separate variables.

\paragraph{CS3: Capability removal.} Removing one capability at a time reduces sufficiency,
recall or both, and turns detections into abstentions on missing evidence
(Table~\ref{tab:cs3}). Removing unit semantics leaves task sufficiency unchanged in this
budget yet causes the largest loss of fault recall: the effect of a capability depends on
whether it is needed to answer or to check.

\begin{table}[t]
\centering\footnotesize
\setlength{\tabcolsep}{3pt}
\caption{CS4: equally informed procedural and SHACL validators on identical engineering facts.}
\label{tab:cs4}
\begin{tabular}{@{}lrrr@{}}
\toprule
Validator & Faults caught & False alarms & Diagnosis \\
\midrule
Procedural & 84/84 & 0/53 & 100\% \\
SHACL & 84/84 & 0/53 & 100\% \\
\bottomrule
\end{tabular}
\end{table}


\paragraph{CS4: Representation versus information.} A procedural validator and a SHACL
validator given identical engineering facts produce \CsDisagreements{} raw-family
disagreements (Table~\ref{tab:cs4}). The measured benefit lies in explicit engineering
evidence, not in the validation technology; reuse, auditability and maintenance effort
remain to be measured.

\begin{table}[t]
\centering\footnotesize
\setlength{\tabcolsep}{3pt}
\caption{CS5: matched facts recovered from the emitted OPC UA NodeSet2 and AAS v3 projections of the same plant ($\neq$: conflicting value). Both projections are generated from one specification, so this measures the emitted profiles, not limits of either standard.}
\label{tab:cs5}
\begin{tabular}{@{}llrrr@{}}
\toprule
Fact & AUF & OPC UA & AAS & Union \\
\midrule
identity & U1 & 73/73 & 73/73 & 73/73 \\
semantic id & U1 & 0/46 & 46/46 & 46/46 \\
nameplate & U1 & 0/11 & 11/11 & 11/11 \\
datatype & U2 & 73/73 & 73/73 & 73/73 \\
access & U2 & 73/73 & 0/73 & 73/73 \\
range & U2 & 73/73 & 0/73 & 73/73 \\
unit & U3 & 73/73 & 73/73 & 73/73 \\
quantity kind & U3 & 53/73 (20$\neq$) & 0/73 & 53/73 (20$\neq$) \\
membership & U4 & 73/73 & 73/73 & 73/73 \\
parent & U4 & 16/16 & 0/16 & 16/16 \\
feeds & U4 & 0/14 & 0/14 & 0/14 \\
state & U5 & 0/9 & 9/9 & 9/9 \\
legal commands & U5 & 0/9 & 0/9 & 0/9 \\
interlock & U6 & 0/4 & 0/4 & 0/4 \\
value & U7 & 0/53 & 53/53 & 53/53 \\
observation time & U7 & 0/53 & 0/53 & 0/53 \\
\bottomrule
\end{tabular}
\end{table}


\paragraph{CS5: OPC~UA versus AAS.} The plant's \opcua{} NodeSet2 and AAS v3 projections
were parsed back into \CsFacts{} matched facts (Table~\ref{tab:cs5}). The \opcua{} profile
recovers \CsOpcuaFacts{} (\CsOpcuaPct{}), including access levels, engineering ranges and
hierarchy; the AAS profile recovers \CsAasFacts{} (\CsAasPct{}), including semantic
identifiers, nameplates, operating state and values; their union reaches \CsUnionFacts{}
(\CsUnionPct{}). \opcua{} \texttt{EUInformation} is attached to the unit, so
\CsQkConflicts{} count-type signals report the unit's quantity kind rather than the
signal's. Neither projection carries flow topology, legal commands, interlocks or
observation timestamps, which C4 requires. The serialisations cost \CsOpcuaTokens{} and
\CsAasTokens{} surrogate tokens. Both standards can express more through companion models,
qualifiers or additional submodels; the result shows that the emitted profiles are
complementary and jointly incomplete, not that either standard is inferior. A second,
derived pilot, an IIH-attributed KUKA robot NodeSet and its loss-aware AAS projection,
retains all \CsKukaRetained{} of \CsKukaFacts{} source facts, including 90 nodes and
277 references; it verifies conversion fidelity only.
